\chapter{原子模型和单电子原子}

\problem{1.1}{氯化钠晶体组成立方点阵，钠和氯原子沿三个轴交错占据位置。已知它们的原子量分别为 22.99 和 35.46，氯化钠的密度为 $2.17\times10^3\ \mathrm{kg\,m^{-3}}$，试估算两相邻离子的间隔。}
% layout
\answerarea

\problem{1.2}{\begin{enumerate}
\item 广播天线以频率为 $1\ \mathrm{MHz}$、功率为 $1\ \mathrm{kW}$ 发射无线电波。求每秒发射的光子数。
\item 太阳垂直入射到地球表面上的辐射率是 $1.94\ \mathrm{cal\,cm^{-2}\,min^{-1}}$，平均波长为 $550\ \mathrm{nm}$。如果直接去看它，设眼球接收光的面积为 $1\ \mathrm{cm^2}$，求每秒内人的眼睛接收到的光子数。这两个数目表明为什么在研究广播辐射和太阳光学时，电磁辐射的量子特性并未直接显示出来。
\end{enumerate}}
% layout
\answerarea
\newpage

\problem{1.3}{已知天空中相当明亮的一等星在地球表面产生大约 $10^{-6}\ \mathrm{lm\,m^{-2}}$ 的光通量，$1\ \mathrm{lm}$ 平均波长为 $556\ \mathrm{nm}$ 的光相当于 $0.016\ \mathrm{W}$，正常人眼只要接收到 100 个光子就有感觉。试估算每秒进入人眼中的光子数，并说明天上某些星星的“眨眼”（包括一等星）是否是由于光的量子性引起的。}
% layout
\answerarea

\problem{1.4}{在密立根的油滴实验中，一个特定的油滴在两块相距为 $5\ \mathrm{mm}$ 的水平极板之间自由下落，速度为 $2.26\times10^{-4}\ \mathrm{m\,s^{-1}}$。在两板上加一电势差 $1600\ \mathrm{V}$ 后，油滴以 $0.90\times10^{-4}\ \mathrm{m\,s^{-1}}$ 速度均匀上升。已知空气的黏滞度为 $1.80\times10^{-5}\ \mathrm{N\,s\,m^{-2}}$，油的密度为 $900\ \mathrm{kg\,m^{-3}}$，试求油滴的半径和它所带的电荷。}
% layout
\answerarea
\newpage

\problem{1.5}{\begin{enumerate}
\item 设有正电荷均匀分布在一半径为 $R$ 的球形区域内，电荷密度为 $\rho$，试证明电荷为 $-e$ 的电子在它内部可以做围绕球心的简谐运动。
\item 若正电荷大小等于电子电荷，$R=1.0\times10^{-10}\ \mathrm{m}$，求作用力常数 $k$ 和电子的振动频率。
\end{enumerate}}
% layout
\answerarea

\problem{1.6}{在 $\alpha$ 粒子散射实验中，若 $\alpha$ 放射源用的是 $^{210}\mathrm{Po}$，它发出的 $\alpha$ 粒子能量为 $5.30\ \mathrm{MeV}$，散射体用 $Z=79$ 的金箔，求：
\begin{enumerate}
\item 散射角为 $90^\circ$ 所对应的瞄准距离；
\item 在这种情况下，$\alpha$ 粒子与金核达到的最短距离；
\item 这种能量下的 $\alpha$ 粒子与金核能达到的最短距离。
\end{enumerate}}
% layout
\answerarea
\newpage

\problem{1.7}{加速器产生的能量为 $1.5\ \mathrm{MeV}$ 质子束垂直入射到厚为 $1\ \mu\mathrm{m}$ 的金箔上，求：
\begin{enumerate}
\item 散射角大于 $90^\circ$ 的质子数占全部入射粒子数的百分比（已知金的 $A=197$，$\rho=1.932\times10^4\ \mathrm{kg\,m^{-3}}$）；
\item 散射到 $29^\circ\!\sim31^\circ$、$89^\circ\!\sim91^\circ$ 和 $149^\circ\!\sim151^\circ$ 间隔内质子数的比例；
\item 设束流为 $10\ \mathrm{nA}$，10 分钟内散射到 $149^\circ\!\sim151^\circ$ 内的质子数；
\item 在 $30^\circ$ 角方向质子的卢瑟福散射微分截面。
\end{enumerate}}
% layout
\answerarea
\newpage

\problem{1.8}{计算氢原子基态中电子绕核运动的轨道半径、频率、周期、线速度、线动量、角速度、角动量、加速度、动能、势能和总能。}
% layout
\answerarea

\problem{1.9}{在氢原子中，求：
\begin{enumerate}
\item 电子在 $n=1$ 轨道上运动时相应的电流值大小；
\item $n=1$ 轨道中心处的磁场强度。
\end{enumerate}}
% layout
\answerarea
\newpage

\problem{1.10}{试求氢原子的电子与核之间的库仑引力和万有引力以及它们的比值，从中看出把万有引力略去不计是否合理。}
% layout
\answerarea

\problem{1.11}{分别计算氢的赖曼系、巴耳末系和帕邢系中最短和最长的波长。这三个线系各属于哪个电磁波段？}
% layout
\answerarea
\newpage

\problem{1.12}{求氦离子 $\mathrm{He}^{+}$ 的电离能、第一和第二激发能、赖曼系第一条谱线的波长、巴耳末系限波长和第一玻尔轨道半径。}
% layout
\answerarea

\problem{1.13}{从含有氢和氦的放电管内得到的光谱中，发现有一条线离氢的 $\mathrm{H}_{\alpha}$ 线（$656.279\ \mathrm{nm}$）的距离为 $2.674\times10^{-10}\ \mathrm{m}$，这被归于一次电离的氦原子的跃迁。
\begin{enumerate}
\item 找出 $\mathrm{He}^{+}$ 的这一跃迁中涉及的能级的主量子数；
\item 计算离子 $\mathrm{He}^{+}$ 的里德伯常数（已知 $R_{\mathrm H}=1.096\,775\,8\times10^7\ \mathrm{m^{-1}}$，氦核质量为 $3726.358\ \mathrm{MeV}$）；
\item 假设核为无限重的，计算 $\mathrm{He}^{+}$ 的里德伯常数。
\end{enumerate}}
% layout
\answerarea
\newpage

\problem{1.14}{某一气体放电管中装有 $^{1}\mathrm H$、$^{2}\mathrm H$、$^{3}\mathrm{He}^{+}$、$^{4}\mathrm{He}^{+}$、$^{6}\mathrm{Li}^{2+}$ 和 $^{7}\mathrm{Li}^{2+}$ 原子或离子。
\begin{enumerate}
\item 计算它们的电离能量；
\item 当放电管所加的电压从零逐渐增大时，最先出现哪一条谱线？
\end{enumerate}}
% layout
\answerarea

\problem{1.15}{用 $12.9\ \mathrm{eV}$ 的电子去激发基态氢原子。
\begin{enumerate}
\item 求受激发的氢原子向低能级跃迁时发出的光谱线；
\item 如果这个氢原子最初是静止的，计算当它从 $n=3$ 态直接跃迁到 $n=1$ 态时的反冲能量和速度。
\end{enumerate}}
% layout
\answerarea
\newpage

\problem{1.16}{设一 $\mu^-$ 子代替氢原子中一个电子而形成中性氢 $\mu$ 原子，试求它的基态结合能和从 $n=3$ 态到 $n=2$ 态跃迁发出的光子能量。如果是中性锂 $\mu$ 原子，它的化学性质最类似于哪种化学元素？基态结合能呢？}
% layout
\answerarea

\problem{1.17}{\begin{enumerate}
\item 用最简单的氢原子理论粗略估算 $Z=50$ 的 $\pi$ 原子的最低两个能级所对应的能量和轨道半径、电离能和第一激发能量；
\item 将这些结果与氢原子的能量和轨道半径进行比较。已知 $m_{\pi}=273m_e$。做这种比较是有意义的，由此可以得到一些什么结论？
\end{enumerate}}
% layout
\answerarea
\newpage

\problem{1.18}{\begin{enumerate}
\item 试求钠原子被激发到 $n=100$ 的里德伯原子态的原子半径、电离能和第一激发能；
\item 将所求的结果与氢原子的 $n=100$ 的里德伯原子态所对应的量作一比较。
\end{enumerate}}
% layout
\answerarea

\newpage

